A apicultura no semiárido cearense, embora vantajosa, enfrenta barreiras de competitividade que o associativismo, baseado na Economia Solidária, ajuda a superar. Contudo, a delegação de funções a uma diretoria gera assimetria informacional entre os associados, o que pode comprometer a confiança e a coesão do grupo. Este trabalho tem como objetivo propor uma aplicação web para automatizar processos internos da Associação de Apicultores de Poranga, fortalecendo a transparência da gestão. Dada a baixa familiaridade dos associados com ferramentas digitais, os requisitos serão levantados por meio de Prototipação Exploratória, conduzida em ciclos iterativos com stakeholders representativos dos principais perfis de usuário, com o protótipo construído segundo a abordagem Mobile First. O desenvolvimento seguirá o ciclo Test-Driven Development (TDD), com testes automatizados (unitários, de integração e de regressão) garantindo a confiabilidade do sistema ao longo de toda a evolução do código, complementados por uma avaliação de usabilidade conduzida com associados não envolvidos na etapa de elicitação. Espera-se que, ao final do processo, a aplicação seja capaz de automatizar os processos internos da associação com confiabilidade, disponibilizando informações de forma transparente e acessível a todos os associados, de modo a reduzir a assimetria informacional, fortalecer a confiança mútua e sustentar a autogestão coletiva da organização.
 
\textbf{Palavras-chave:} Apicultura. Economia Solidária. Associativismo. Engenharia de Software. Test-Driven Development.